\documentclass[11pt,a4paper]{article}

\newcommand{\CourseCode}{CSAI2017P: Applied Machine Learning (Lab)}
\newcommand{\AssignmentName}{LA03 -- Lab Assignment 3}

% Shared settings for Applied Machine Learning lab assignments and marking schemes.
% Layered on the house notes preamble; keep free of \documentclass.
%
% Callers must define \CourseCode, \AssignmentName before \begin{document}.

% Depth-agnostic locate of the house notes preamble: briefs compile from
% public/lab-assignments/LAxx/latex/ and marking schemes from
% private/solutions/lab-assignments/LAxx/latex/, which sit at different depths.
\IfFileExists{../../../../public/_shared/notes-common.tex}{%
  % Shared notes settings for Applied Machine Learning lectures (UPES).
% Keep this file free of \documentclass to allow reuse via \input.

\usepackage[utf8]{inputenc}
\usepackage[T1]{fontenc}
\usepackage{lmodern}          % scalable T1 fonts; without this pdflatex falls back
\input{glyphtounicode}        % to bitmapped EC Type 3 and the PDFs are neither
\pdfgentounicode=1            % searchable nor screen-reader friendly
\usepackage[a4paper,margin=1in]{geometry}
\usepackage{hyperref}
\usepackage{xurl}
\usepackage{booktabs}
\usepackage{amsmath}
\usepackage{amssymb}
\usepackage{listings}
\usepackage{xcolor}
\usepackage{enumitem}
\usepackage{parskip}

% Code listings (Python-oriented for this subject).
\lstset{
  language=Python,
  basicstyle=\ttfamily\small,
  keywordstyle=\color{blue},
  commentstyle=\color{gray},
  stringstyle=\color{teal},
  showstringspaces=false,
  columns=fullflexible,
  frame=single,
  framerule=0.2pt,
  breaklines=true
}

\setlist[itemize]{topsep=0.3em,itemsep=0.2em}
\setlist[enumerate]{topsep=0.3em,itemsep=0.2em}

% Simple per-section source line.
\newcommand{\notesources}[1]{\par\noindent{\small\color{gray}\textit{Sources:} \sloppy #1\par}}

% Unit-wise source strings for this subject (used by slides + notes).
% Path is relative to the lecture `latex/` folder (the compilation CWD).
\IfFileExists{../../../_shared/aml-sources.tex}{%
  % Source strings for Applied Machine Learning (CSAI2017P) — used by slides + notes.
% Prescribed texts and reference books are taken from the course syllabus.
% Support texts are the author-hosted open-access editions:
%   statlearning.com            (ISL with Python)
%   probml.github.io/pml-book   (Murphy, Probabilistic Machine Learning)
%   d2l.ai                      (Dive into Deep Learning)
%   mml-book.github.io          (Mathematics for Machine Learning)

\newcommand{\AMLRepoURL}{https://github.com/tali7c/Applied-Machine-Learning}

% --- prescribed -------------------------------------------------------------
\newcommand{\AMLPrescribed}{%
M\"uller and Guido, \textit{Introduction to Machine Learning with Python}, O'Reilly, 2016.;%
Bishop, \textit{Pattern Recognition and Machine Learning}, Springer, 2006.;%
Murphy, \textit{Machine Learning: A Probabilistic Perspective}, MIT Press, 2012.;%
Han, Kamber and Pei, \textit{Data Mining: Concepts and Techniques} (3e), Morgan Kaufmann, 2012.%
}

% --- unit-wise ---------------------------------------------------------------
\newcommand{\AMLUnitISources}{%
M\"uller and Guido, \textit{Introduction to Machine Learning with Python}, Ch. 1.;%
Bishop, \textit{Pattern Recognition and Machine Learning}, Ch. 1.;%
Murphy, \textit{Probabilistic Machine Learning: An Introduction}, MIT Press, 2022, Ch. 1.;%
Mitchell, \textit{Machine Learning}, McGraw Hill, 1997, Ch. 1.;%
Scikit-learn documentation, accessed 2026-08-04.%
}

\newcommand{\AMLUnitIISources}{%
Bishop, \textit{Pattern Recognition and Machine Learning}, \S1.5, \S4.3, \S7.1.;%
Zhang, Lipton, Li and Smola, \textit{Dive into Deep Learning}, \S3.4.;%
Shalev-Shwartz and Ben-David, \textit{Understanding Machine Learning}, Ch. 15.;%
Murphy, \textit{Probabilistic Machine Learning: An Introduction}, Ch. 5.%
}

\newcommand{\AMLUnitIIISources}{%
Zhang, Lipton, Li and Smola, \textit{Dive into Deep Learning}, Ch. 12 (Optimization Algorithms).;%
Deisenroth, Faisal and Ong, \textit{Mathematics for Machine Learning}, Ch. 7.;%
Kingma and Ba, ``Adam: A Method for Stochastic Optimization,'' ICLR 2015.;%
Ruder, ``An overview of gradient descent optimization algorithms,'' arXiv:1609.04747, 2016.%
}

\newcommand{\AMLUnitIVSources}{%
Han, Kamber and Pei, \textit{Data Mining: Concepts and Techniques} (3e), Ch. 3.;%
M\"uller and Guido, \textit{Introduction to Machine Learning with Python}, Ch. 3--4.;%
James, Witten, Hastie, Tibshirani and Taylor, \textit{An Introduction to Statistical Learning with Python}, Ch. 5--6, 12.;%
Scikit-learn documentation (preprocessing, model selection), accessed 2026-08-04.%
}

\newcommand{\AMLUnitVSources}{%
James et al., \textit{An Introduction to Statistical Learning with Python}, Ch. 3--6.;%
Bishop, \textit{Pattern Recognition and Machine Learning}, Ch. 3.;%
Deisenroth, Faisal and Ong, \textit{Mathematics for Machine Learning}, Ch. 9.;%
M\"uller and Guido, \textit{Introduction to Machine Learning with Python}, \S2.3.%
}

\newcommand{\AMLUnitVISources}{%
Han, Kamber and Pei, \textit{Data Mining: Concepts and Techniques} (3e), Ch. 8--9.;%
James et al., \textit{An Introduction to Statistical Learning with Python}, Ch. 8--9.;%
Bishop, \textit{Pattern Recognition and Machine Learning}, Ch. 4--5, 7, 14.;%
Zhang et al., \textit{Dive into Deep Learning}, Ch. 5.%
}

\newcommand{\AMLUnitVIISources}{%
Han, Kamber and Pei, \textit{Data Mining: Concepts and Techniques} (3e), Ch. 10.;%
James et al., \textit{An Introduction to Statistical Learning with Python}, Ch. 12.;%
M\"uller and Guido, \textit{Introduction to Machine Learning with Python}, \S3.5.;%
Ester, Kriegel, Sander and Xu, ``A density-based algorithm for discovering clusters (DBSCAN),'' KDD 1996.%
}

\newcommand{\AMLLabSources}{%
M\"uller and Guido, \textit{Introduction to Machine Learning with Python}, companion notebooks.;%
McKinney, \textit{Python for Data Analysis} (3e), O'Reilly, 2022.;%
Scikit-learn user guide, accessed 2026-08-04.;%
Matplotlib and Seaborn documentation, accessed 2026-08-04.%
}
%
}{%
  % Faculty copies live under `private/` so their compile CWD is different.
  \IfFileExists{../../../../public/_shared/aml-sources.tex}{%
    % Source strings for Applied Machine Learning (CSAI2017P) — used by slides + notes.
% Prescribed texts and reference books are taken from the course syllabus.
% Support texts are the author-hosted open-access editions:
%   statlearning.com            (ISL with Python)
%   probml.github.io/pml-book   (Murphy, Probabilistic Machine Learning)
%   d2l.ai                      (Dive into Deep Learning)
%   mml-book.github.io          (Mathematics for Machine Learning)

\newcommand{\AMLRepoURL}{https://github.com/tali7c/Applied-Machine-Learning}

% --- prescribed -------------------------------------------------------------
\newcommand{\AMLPrescribed}{%
M\"uller and Guido, \textit{Introduction to Machine Learning with Python}, O'Reilly, 2016.;%
Bishop, \textit{Pattern Recognition and Machine Learning}, Springer, 2006.;%
Murphy, \textit{Machine Learning: A Probabilistic Perspective}, MIT Press, 2012.;%
Han, Kamber and Pei, \textit{Data Mining: Concepts and Techniques} (3e), Morgan Kaufmann, 2012.%
}

% --- unit-wise ---------------------------------------------------------------
\newcommand{\AMLUnitISources}{%
M\"uller and Guido, \textit{Introduction to Machine Learning with Python}, Ch. 1.;%
Bishop, \textit{Pattern Recognition and Machine Learning}, Ch. 1.;%
Murphy, \textit{Probabilistic Machine Learning: An Introduction}, MIT Press, 2022, Ch. 1.;%
Mitchell, \textit{Machine Learning}, McGraw Hill, 1997, Ch. 1.;%
Scikit-learn documentation, accessed 2026-08-04.%
}

\newcommand{\AMLUnitIISources}{%
Bishop, \textit{Pattern Recognition and Machine Learning}, \S1.5, \S4.3, \S7.1.;%
Zhang, Lipton, Li and Smola, \textit{Dive into Deep Learning}, \S3.4.;%
Shalev-Shwartz and Ben-David, \textit{Understanding Machine Learning}, Ch. 15.;%
Murphy, \textit{Probabilistic Machine Learning: An Introduction}, Ch. 5.%
}

\newcommand{\AMLUnitIIISources}{%
Zhang, Lipton, Li and Smola, \textit{Dive into Deep Learning}, Ch. 12 (Optimization Algorithms).;%
Deisenroth, Faisal and Ong, \textit{Mathematics for Machine Learning}, Ch. 7.;%
Kingma and Ba, ``Adam: A Method for Stochastic Optimization,'' ICLR 2015.;%
Ruder, ``An overview of gradient descent optimization algorithms,'' arXiv:1609.04747, 2016.%
}

\newcommand{\AMLUnitIVSources}{%
Han, Kamber and Pei, \textit{Data Mining: Concepts and Techniques} (3e), Ch. 3.;%
M\"uller and Guido, \textit{Introduction to Machine Learning with Python}, Ch. 3--4.;%
James, Witten, Hastie, Tibshirani and Taylor, \textit{An Introduction to Statistical Learning with Python}, Ch. 5--6, 12.;%
Scikit-learn documentation (preprocessing, model selection), accessed 2026-08-04.%
}

\newcommand{\AMLUnitVSources}{%
James et al., \textit{An Introduction to Statistical Learning with Python}, Ch. 3--6.;%
Bishop, \textit{Pattern Recognition and Machine Learning}, Ch. 3.;%
Deisenroth, Faisal and Ong, \textit{Mathematics for Machine Learning}, Ch. 9.;%
M\"uller and Guido, \textit{Introduction to Machine Learning with Python}, \S2.3.%
}

\newcommand{\AMLUnitVISources}{%
Han, Kamber and Pei, \textit{Data Mining: Concepts and Techniques} (3e), Ch. 8--9.;%
James et al., \textit{An Introduction to Statistical Learning with Python}, Ch. 8--9.;%
Bishop, \textit{Pattern Recognition and Machine Learning}, Ch. 4--5, 7, 14.;%
Zhang et al., \textit{Dive into Deep Learning}, Ch. 5.%
}

\newcommand{\AMLUnitVIISources}{%
Han, Kamber and Pei, \textit{Data Mining: Concepts and Techniques} (3e), Ch. 10.;%
James et al., \textit{An Introduction to Statistical Learning with Python}, Ch. 12.;%
M\"uller and Guido, \textit{Introduction to Machine Learning with Python}, \S3.5.;%
Ester, Kriegel, Sander and Xu, ``A density-based algorithm for discovering clusters (DBSCAN),'' KDD 1996.%
}

\newcommand{\AMLLabSources}{%
M\"uller and Guido, \textit{Introduction to Machine Learning with Python}, companion notebooks.;%
McKinney, \textit{Python for Data Analysis} (3e), O'Reilly, 2022.;%
Scikit-learn user guide, accessed 2026-08-04.;%
Matplotlib and Seaborn documentation, accessed 2026-08-04.%
}
%
  }{%
    % Question papers live at `private/<family>/<instrument>/`.
    \IfFileExists{../../../public/_shared/aml-sources.tex}{%
      % Source strings for Applied Machine Learning (CSAI2017P) — used by slides + notes.
% Prescribed texts and reference books are taken from the course syllabus.
% Support texts are the author-hosted open-access editions:
%   statlearning.com            (ISL with Python)
%   probml.github.io/pml-book   (Murphy, Probabilistic Machine Learning)
%   d2l.ai                      (Dive into Deep Learning)
%   mml-book.github.io          (Mathematics for Machine Learning)

\newcommand{\AMLRepoURL}{https://github.com/tali7c/Applied-Machine-Learning}

% --- prescribed -------------------------------------------------------------
\newcommand{\AMLPrescribed}{%
M\"uller and Guido, \textit{Introduction to Machine Learning with Python}, O'Reilly, 2016.;%
Bishop, \textit{Pattern Recognition and Machine Learning}, Springer, 2006.;%
Murphy, \textit{Machine Learning: A Probabilistic Perspective}, MIT Press, 2012.;%
Han, Kamber and Pei, \textit{Data Mining: Concepts and Techniques} (3e), Morgan Kaufmann, 2012.%
}

% --- unit-wise ---------------------------------------------------------------
\newcommand{\AMLUnitISources}{%
M\"uller and Guido, \textit{Introduction to Machine Learning with Python}, Ch. 1.;%
Bishop, \textit{Pattern Recognition and Machine Learning}, Ch. 1.;%
Murphy, \textit{Probabilistic Machine Learning: An Introduction}, MIT Press, 2022, Ch. 1.;%
Mitchell, \textit{Machine Learning}, McGraw Hill, 1997, Ch. 1.;%
Scikit-learn documentation, accessed 2026-08-04.%
}

\newcommand{\AMLUnitIISources}{%
Bishop, \textit{Pattern Recognition and Machine Learning}, \S1.5, \S4.3, \S7.1.;%
Zhang, Lipton, Li and Smola, \textit{Dive into Deep Learning}, \S3.4.;%
Shalev-Shwartz and Ben-David, \textit{Understanding Machine Learning}, Ch. 15.;%
Murphy, \textit{Probabilistic Machine Learning: An Introduction}, Ch. 5.%
}

\newcommand{\AMLUnitIIISources}{%
Zhang, Lipton, Li and Smola, \textit{Dive into Deep Learning}, Ch. 12 (Optimization Algorithms).;%
Deisenroth, Faisal and Ong, \textit{Mathematics for Machine Learning}, Ch. 7.;%
Kingma and Ba, ``Adam: A Method for Stochastic Optimization,'' ICLR 2015.;%
Ruder, ``An overview of gradient descent optimization algorithms,'' arXiv:1609.04747, 2016.%
}

\newcommand{\AMLUnitIVSources}{%
Han, Kamber and Pei, \textit{Data Mining: Concepts and Techniques} (3e), Ch. 3.;%
M\"uller and Guido, \textit{Introduction to Machine Learning with Python}, Ch. 3--4.;%
James, Witten, Hastie, Tibshirani and Taylor, \textit{An Introduction to Statistical Learning with Python}, Ch. 5--6, 12.;%
Scikit-learn documentation (preprocessing, model selection), accessed 2026-08-04.%
}

\newcommand{\AMLUnitVSources}{%
James et al., \textit{An Introduction to Statistical Learning with Python}, Ch. 3--6.;%
Bishop, \textit{Pattern Recognition and Machine Learning}, Ch. 3.;%
Deisenroth, Faisal and Ong, \textit{Mathematics for Machine Learning}, Ch. 9.;%
M\"uller and Guido, \textit{Introduction to Machine Learning with Python}, \S2.3.%
}

\newcommand{\AMLUnitVISources}{%
Han, Kamber and Pei, \textit{Data Mining: Concepts and Techniques} (3e), Ch. 8--9.;%
James et al., \textit{An Introduction to Statistical Learning with Python}, Ch. 8--9.;%
Bishop, \textit{Pattern Recognition and Machine Learning}, Ch. 4--5, 7, 14.;%
Zhang et al., \textit{Dive into Deep Learning}, Ch. 5.%
}

\newcommand{\AMLUnitVIISources}{%
Han, Kamber and Pei, \textit{Data Mining: Concepts and Techniques} (3e), Ch. 10.;%
James et al., \textit{An Introduction to Statistical Learning with Python}, Ch. 12.;%
M\"uller and Guido, \textit{Introduction to Machine Learning with Python}, \S3.5.;%
Ester, Kriegel, Sander and Xu, ``A density-based algorithm for discovering clusters (DBSCAN),'' KDD 1996.%
}

\newcommand{\AMLLabSources}{%
M\"uller and Guido, \textit{Introduction to Machine Learning with Python}, companion notebooks.;%
McKinney, \textit{Python for Data Analysis} (3e), O'Reilly, 2022.;%
Scikit-learn user guide, accessed 2026-08-04.;%
Matplotlib and Seaborn documentation, accessed 2026-08-04.%
}
%
    }{%
      % Marking schemes live at `private/solutions/<family>/<id>/latex/`.
      \IfFileExists{../../../../../public/_shared/aml-sources.tex}{%
        % Source strings for Applied Machine Learning (CSAI2017P) — used by slides + notes.
% Prescribed texts and reference books are taken from the course syllabus.
% Support texts are the author-hosted open-access editions:
%   statlearning.com            (ISL with Python)
%   probml.github.io/pml-book   (Murphy, Probabilistic Machine Learning)
%   d2l.ai                      (Dive into Deep Learning)
%   mml-book.github.io          (Mathematics for Machine Learning)

\newcommand{\AMLRepoURL}{https://github.com/tali7c/Applied-Machine-Learning}

% --- prescribed -------------------------------------------------------------
\newcommand{\AMLPrescribed}{%
M\"uller and Guido, \textit{Introduction to Machine Learning with Python}, O'Reilly, 2016.;%
Bishop, \textit{Pattern Recognition and Machine Learning}, Springer, 2006.;%
Murphy, \textit{Machine Learning: A Probabilistic Perspective}, MIT Press, 2012.;%
Han, Kamber and Pei, \textit{Data Mining: Concepts and Techniques} (3e), Morgan Kaufmann, 2012.%
}

% --- unit-wise ---------------------------------------------------------------
\newcommand{\AMLUnitISources}{%
M\"uller and Guido, \textit{Introduction to Machine Learning with Python}, Ch. 1.;%
Bishop, \textit{Pattern Recognition and Machine Learning}, Ch. 1.;%
Murphy, \textit{Probabilistic Machine Learning: An Introduction}, MIT Press, 2022, Ch. 1.;%
Mitchell, \textit{Machine Learning}, McGraw Hill, 1997, Ch. 1.;%
Scikit-learn documentation, accessed 2026-08-04.%
}

\newcommand{\AMLUnitIISources}{%
Bishop, \textit{Pattern Recognition and Machine Learning}, \S1.5, \S4.3, \S7.1.;%
Zhang, Lipton, Li and Smola, \textit{Dive into Deep Learning}, \S3.4.;%
Shalev-Shwartz and Ben-David, \textit{Understanding Machine Learning}, Ch. 15.;%
Murphy, \textit{Probabilistic Machine Learning: An Introduction}, Ch. 5.%
}

\newcommand{\AMLUnitIIISources}{%
Zhang, Lipton, Li and Smola, \textit{Dive into Deep Learning}, Ch. 12 (Optimization Algorithms).;%
Deisenroth, Faisal and Ong, \textit{Mathematics for Machine Learning}, Ch. 7.;%
Kingma and Ba, ``Adam: A Method for Stochastic Optimization,'' ICLR 2015.;%
Ruder, ``An overview of gradient descent optimization algorithms,'' arXiv:1609.04747, 2016.%
}

\newcommand{\AMLUnitIVSources}{%
Han, Kamber and Pei, \textit{Data Mining: Concepts and Techniques} (3e), Ch. 3.;%
M\"uller and Guido, \textit{Introduction to Machine Learning with Python}, Ch. 3--4.;%
James, Witten, Hastie, Tibshirani and Taylor, \textit{An Introduction to Statistical Learning with Python}, Ch. 5--6, 12.;%
Scikit-learn documentation (preprocessing, model selection), accessed 2026-08-04.%
}

\newcommand{\AMLUnitVSources}{%
James et al., \textit{An Introduction to Statistical Learning with Python}, Ch. 3--6.;%
Bishop, \textit{Pattern Recognition and Machine Learning}, Ch. 3.;%
Deisenroth, Faisal and Ong, \textit{Mathematics for Machine Learning}, Ch. 9.;%
M\"uller and Guido, \textit{Introduction to Machine Learning with Python}, \S2.3.%
}

\newcommand{\AMLUnitVISources}{%
Han, Kamber and Pei, \textit{Data Mining: Concepts and Techniques} (3e), Ch. 8--9.;%
James et al., \textit{An Introduction to Statistical Learning with Python}, Ch. 8--9.;%
Bishop, \textit{Pattern Recognition and Machine Learning}, Ch. 4--5, 7, 14.;%
Zhang et al., \textit{Dive into Deep Learning}, Ch. 5.%
}

\newcommand{\AMLUnitVIISources}{%
Han, Kamber and Pei, \textit{Data Mining: Concepts and Techniques} (3e), Ch. 10.;%
James et al., \textit{An Introduction to Statistical Learning with Python}, Ch. 12.;%
M\"uller and Guido, \textit{Introduction to Machine Learning with Python}, \S3.5.;%
Ester, Kriegel, Sander and Xu, ``A density-based algorithm for discovering clusters (DBSCAN),'' KDD 1996.%
}

\newcommand{\AMLLabSources}{%
M\"uller and Guido, \textit{Introduction to Machine Learning with Python}, companion notebooks.;%
McKinney, \textit{Python for Data Analysis} (3e), O'Reilly, 2022.;%
Scikit-learn user guide, accessed 2026-08-04.;%
Matplotlib and Seaborn documentation, accessed 2026-08-04.%
}
%
      }{%
        % Source strings for Applied Machine Learning (CSAI2017P) — used by slides + notes.
% Prescribed texts and reference books are taken from the course syllabus.
% Support texts are the author-hosted open-access editions:
%   statlearning.com            (ISL with Python)
%   probml.github.io/pml-book   (Murphy, Probabilistic Machine Learning)
%   d2l.ai                      (Dive into Deep Learning)
%   mml-book.github.io          (Mathematics for Machine Learning)

\newcommand{\AMLRepoURL}{https://github.com/tali7c/Applied-Machine-Learning}

% --- prescribed -------------------------------------------------------------
\newcommand{\AMLPrescribed}{%
M\"uller and Guido, \textit{Introduction to Machine Learning with Python}, O'Reilly, 2016.;%
Bishop, \textit{Pattern Recognition and Machine Learning}, Springer, 2006.;%
Murphy, \textit{Machine Learning: A Probabilistic Perspective}, MIT Press, 2012.;%
Han, Kamber and Pei, \textit{Data Mining: Concepts and Techniques} (3e), Morgan Kaufmann, 2012.%
}

% --- unit-wise ---------------------------------------------------------------
\newcommand{\AMLUnitISources}{%
M\"uller and Guido, \textit{Introduction to Machine Learning with Python}, Ch. 1.;%
Bishop, \textit{Pattern Recognition and Machine Learning}, Ch. 1.;%
Murphy, \textit{Probabilistic Machine Learning: An Introduction}, MIT Press, 2022, Ch. 1.;%
Mitchell, \textit{Machine Learning}, McGraw Hill, 1997, Ch. 1.;%
Scikit-learn documentation, accessed 2026-08-04.%
}

\newcommand{\AMLUnitIISources}{%
Bishop, \textit{Pattern Recognition and Machine Learning}, \S1.5, \S4.3, \S7.1.;%
Zhang, Lipton, Li and Smola, \textit{Dive into Deep Learning}, \S3.4.;%
Shalev-Shwartz and Ben-David, \textit{Understanding Machine Learning}, Ch. 15.;%
Murphy, \textit{Probabilistic Machine Learning: An Introduction}, Ch. 5.%
}

\newcommand{\AMLUnitIIISources}{%
Zhang, Lipton, Li and Smola, \textit{Dive into Deep Learning}, Ch. 12 (Optimization Algorithms).;%
Deisenroth, Faisal and Ong, \textit{Mathematics for Machine Learning}, Ch. 7.;%
Kingma and Ba, ``Adam: A Method for Stochastic Optimization,'' ICLR 2015.;%
Ruder, ``An overview of gradient descent optimization algorithms,'' arXiv:1609.04747, 2016.%
}

\newcommand{\AMLUnitIVSources}{%
Han, Kamber and Pei, \textit{Data Mining: Concepts and Techniques} (3e), Ch. 3.;%
M\"uller and Guido, \textit{Introduction to Machine Learning with Python}, Ch. 3--4.;%
James, Witten, Hastie, Tibshirani and Taylor, \textit{An Introduction to Statistical Learning with Python}, Ch. 5--6, 12.;%
Scikit-learn documentation (preprocessing, model selection), accessed 2026-08-04.%
}

\newcommand{\AMLUnitVSources}{%
James et al., \textit{An Introduction to Statistical Learning with Python}, Ch. 3--6.;%
Bishop, \textit{Pattern Recognition and Machine Learning}, Ch. 3.;%
Deisenroth, Faisal and Ong, \textit{Mathematics for Machine Learning}, Ch. 9.;%
M\"uller and Guido, \textit{Introduction to Machine Learning with Python}, \S2.3.%
}

\newcommand{\AMLUnitVISources}{%
Han, Kamber and Pei, \textit{Data Mining: Concepts and Techniques} (3e), Ch. 8--9.;%
James et al., \textit{An Introduction to Statistical Learning with Python}, Ch. 8--9.;%
Bishop, \textit{Pattern Recognition and Machine Learning}, Ch. 4--5, 7, 14.;%
Zhang et al., \textit{Dive into Deep Learning}, Ch. 5.%
}

\newcommand{\AMLUnitVIISources}{%
Han, Kamber and Pei, \textit{Data Mining: Concepts and Techniques} (3e), Ch. 10.;%
James et al., \textit{An Introduction to Statistical Learning with Python}, Ch. 12.;%
M\"uller and Guido, \textit{Introduction to Machine Learning with Python}, \S3.5.;%
Ester, Kriegel, Sander and Xu, ``A density-based algorithm for discovering clusters (DBSCAN),'' KDD 1996.%
}

\newcommand{\AMLLabSources}{%
M\"uller and Guido, \textit{Introduction to Machine Learning with Python}, companion notebooks.;%
McKinney, \textit{Python for Data Analysis} (3e), O'Reilly, 2022.;%
Scikit-learn user guide, accessed 2026-08-04.;%
Matplotlib and Seaborn documentation, accessed 2026-08-04.%
}
%
      }%
    }%
  }%
}
%
}{%
  \IfFileExists{../../../../../public/_shared/notes-common.tex}{%
    % Shared notes settings for Applied Machine Learning lectures (UPES).
% Keep this file free of \documentclass to allow reuse via \input.

\usepackage[utf8]{inputenc}
\usepackage[T1]{fontenc}
\usepackage{lmodern}          % scalable T1 fonts; without this pdflatex falls back
\input{glyphtounicode}        % to bitmapped EC Type 3 and the PDFs are neither
\pdfgentounicode=1            % searchable nor screen-reader friendly
\usepackage[a4paper,margin=1in]{geometry}
\usepackage{hyperref}
\usepackage{xurl}
\usepackage{booktabs}
\usepackage{amsmath}
\usepackage{amssymb}
\usepackage{listings}
\usepackage{xcolor}
\usepackage{enumitem}
\usepackage{parskip}

% Code listings (Python-oriented for this subject).
\lstset{
  language=Python,
  basicstyle=\ttfamily\small,
  keywordstyle=\color{blue},
  commentstyle=\color{gray},
  stringstyle=\color{teal},
  showstringspaces=false,
  columns=fullflexible,
  frame=single,
  framerule=0.2pt,
  breaklines=true
}

\setlist[itemize]{topsep=0.3em,itemsep=0.2em}
\setlist[enumerate]{topsep=0.3em,itemsep=0.2em}

% Simple per-section source line.
\newcommand{\notesources}[1]{\par\noindent{\small\color{gray}\textit{Sources:} \sloppy #1\par}}

% Unit-wise source strings for this subject (used by slides + notes).
% Path is relative to the lecture `latex/` folder (the compilation CWD).
\IfFileExists{../../../_shared/aml-sources.tex}{%
  % Source strings for Applied Machine Learning (CSAI2017P) — used by slides + notes.
% Prescribed texts and reference books are taken from the course syllabus.
% Support texts are the author-hosted open-access editions:
%   statlearning.com            (ISL with Python)
%   probml.github.io/pml-book   (Murphy, Probabilistic Machine Learning)
%   d2l.ai                      (Dive into Deep Learning)
%   mml-book.github.io          (Mathematics for Machine Learning)

\newcommand{\AMLRepoURL}{https://github.com/tali7c/Applied-Machine-Learning}

% --- prescribed -------------------------------------------------------------
\newcommand{\AMLPrescribed}{%
M\"uller and Guido, \textit{Introduction to Machine Learning with Python}, O'Reilly, 2016.;%
Bishop, \textit{Pattern Recognition and Machine Learning}, Springer, 2006.;%
Murphy, \textit{Machine Learning: A Probabilistic Perspective}, MIT Press, 2012.;%
Han, Kamber and Pei, \textit{Data Mining: Concepts and Techniques} (3e), Morgan Kaufmann, 2012.%
}

% --- unit-wise ---------------------------------------------------------------
\newcommand{\AMLUnitISources}{%
M\"uller and Guido, \textit{Introduction to Machine Learning with Python}, Ch. 1.;%
Bishop, \textit{Pattern Recognition and Machine Learning}, Ch. 1.;%
Murphy, \textit{Probabilistic Machine Learning: An Introduction}, MIT Press, 2022, Ch. 1.;%
Mitchell, \textit{Machine Learning}, McGraw Hill, 1997, Ch. 1.;%
Scikit-learn documentation, accessed 2026-08-04.%
}

\newcommand{\AMLUnitIISources}{%
Bishop, \textit{Pattern Recognition and Machine Learning}, \S1.5, \S4.3, \S7.1.;%
Zhang, Lipton, Li and Smola, \textit{Dive into Deep Learning}, \S3.4.;%
Shalev-Shwartz and Ben-David, \textit{Understanding Machine Learning}, Ch. 15.;%
Murphy, \textit{Probabilistic Machine Learning: An Introduction}, Ch. 5.%
}

\newcommand{\AMLUnitIIISources}{%
Zhang, Lipton, Li and Smola, \textit{Dive into Deep Learning}, Ch. 12 (Optimization Algorithms).;%
Deisenroth, Faisal and Ong, \textit{Mathematics for Machine Learning}, Ch. 7.;%
Kingma and Ba, ``Adam: A Method for Stochastic Optimization,'' ICLR 2015.;%
Ruder, ``An overview of gradient descent optimization algorithms,'' arXiv:1609.04747, 2016.%
}

\newcommand{\AMLUnitIVSources}{%
Han, Kamber and Pei, \textit{Data Mining: Concepts and Techniques} (3e), Ch. 3.;%
M\"uller and Guido, \textit{Introduction to Machine Learning with Python}, Ch. 3--4.;%
James, Witten, Hastie, Tibshirani and Taylor, \textit{An Introduction to Statistical Learning with Python}, Ch. 5--6, 12.;%
Scikit-learn documentation (preprocessing, model selection), accessed 2026-08-04.%
}

\newcommand{\AMLUnitVSources}{%
James et al., \textit{An Introduction to Statistical Learning with Python}, Ch. 3--6.;%
Bishop, \textit{Pattern Recognition and Machine Learning}, Ch. 3.;%
Deisenroth, Faisal and Ong, \textit{Mathematics for Machine Learning}, Ch. 9.;%
M\"uller and Guido, \textit{Introduction to Machine Learning with Python}, \S2.3.%
}

\newcommand{\AMLUnitVISources}{%
Han, Kamber and Pei, \textit{Data Mining: Concepts and Techniques} (3e), Ch. 8--9.;%
James et al., \textit{An Introduction to Statistical Learning with Python}, Ch. 8--9.;%
Bishop, \textit{Pattern Recognition and Machine Learning}, Ch. 4--5, 7, 14.;%
Zhang et al., \textit{Dive into Deep Learning}, Ch. 5.%
}

\newcommand{\AMLUnitVIISources}{%
Han, Kamber and Pei, \textit{Data Mining: Concepts and Techniques} (3e), Ch. 10.;%
James et al., \textit{An Introduction to Statistical Learning with Python}, Ch. 12.;%
M\"uller and Guido, \textit{Introduction to Machine Learning with Python}, \S3.5.;%
Ester, Kriegel, Sander and Xu, ``A density-based algorithm for discovering clusters (DBSCAN),'' KDD 1996.%
}

\newcommand{\AMLLabSources}{%
M\"uller and Guido, \textit{Introduction to Machine Learning with Python}, companion notebooks.;%
McKinney, \textit{Python for Data Analysis} (3e), O'Reilly, 2022.;%
Scikit-learn user guide, accessed 2026-08-04.;%
Matplotlib and Seaborn documentation, accessed 2026-08-04.%
}
%
}{%
  % Faculty copies live under `private/` so their compile CWD is different.
  \IfFileExists{../../../../public/_shared/aml-sources.tex}{%
    % Source strings for Applied Machine Learning (CSAI2017P) — used by slides + notes.
% Prescribed texts and reference books are taken from the course syllabus.
% Support texts are the author-hosted open-access editions:
%   statlearning.com            (ISL with Python)
%   probml.github.io/pml-book   (Murphy, Probabilistic Machine Learning)
%   d2l.ai                      (Dive into Deep Learning)
%   mml-book.github.io          (Mathematics for Machine Learning)

\newcommand{\AMLRepoURL}{https://github.com/tali7c/Applied-Machine-Learning}

% --- prescribed -------------------------------------------------------------
\newcommand{\AMLPrescribed}{%
M\"uller and Guido, \textit{Introduction to Machine Learning with Python}, O'Reilly, 2016.;%
Bishop, \textit{Pattern Recognition and Machine Learning}, Springer, 2006.;%
Murphy, \textit{Machine Learning: A Probabilistic Perspective}, MIT Press, 2012.;%
Han, Kamber and Pei, \textit{Data Mining: Concepts and Techniques} (3e), Morgan Kaufmann, 2012.%
}

% --- unit-wise ---------------------------------------------------------------
\newcommand{\AMLUnitISources}{%
M\"uller and Guido, \textit{Introduction to Machine Learning with Python}, Ch. 1.;%
Bishop, \textit{Pattern Recognition and Machine Learning}, Ch. 1.;%
Murphy, \textit{Probabilistic Machine Learning: An Introduction}, MIT Press, 2022, Ch. 1.;%
Mitchell, \textit{Machine Learning}, McGraw Hill, 1997, Ch. 1.;%
Scikit-learn documentation, accessed 2026-08-04.%
}

\newcommand{\AMLUnitIISources}{%
Bishop, \textit{Pattern Recognition and Machine Learning}, \S1.5, \S4.3, \S7.1.;%
Zhang, Lipton, Li and Smola, \textit{Dive into Deep Learning}, \S3.4.;%
Shalev-Shwartz and Ben-David, \textit{Understanding Machine Learning}, Ch. 15.;%
Murphy, \textit{Probabilistic Machine Learning: An Introduction}, Ch. 5.%
}

\newcommand{\AMLUnitIIISources}{%
Zhang, Lipton, Li and Smola, \textit{Dive into Deep Learning}, Ch. 12 (Optimization Algorithms).;%
Deisenroth, Faisal and Ong, \textit{Mathematics for Machine Learning}, Ch. 7.;%
Kingma and Ba, ``Adam: A Method for Stochastic Optimization,'' ICLR 2015.;%
Ruder, ``An overview of gradient descent optimization algorithms,'' arXiv:1609.04747, 2016.%
}

\newcommand{\AMLUnitIVSources}{%
Han, Kamber and Pei, \textit{Data Mining: Concepts and Techniques} (3e), Ch. 3.;%
M\"uller and Guido, \textit{Introduction to Machine Learning with Python}, Ch. 3--4.;%
James, Witten, Hastie, Tibshirani and Taylor, \textit{An Introduction to Statistical Learning with Python}, Ch. 5--6, 12.;%
Scikit-learn documentation (preprocessing, model selection), accessed 2026-08-04.%
}

\newcommand{\AMLUnitVSources}{%
James et al., \textit{An Introduction to Statistical Learning with Python}, Ch. 3--6.;%
Bishop, \textit{Pattern Recognition and Machine Learning}, Ch. 3.;%
Deisenroth, Faisal and Ong, \textit{Mathematics for Machine Learning}, Ch. 9.;%
M\"uller and Guido, \textit{Introduction to Machine Learning with Python}, \S2.3.%
}

\newcommand{\AMLUnitVISources}{%
Han, Kamber and Pei, \textit{Data Mining: Concepts and Techniques} (3e), Ch. 8--9.;%
James et al., \textit{An Introduction to Statistical Learning with Python}, Ch. 8--9.;%
Bishop, \textit{Pattern Recognition and Machine Learning}, Ch. 4--5, 7, 14.;%
Zhang et al., \textit{Dive into Deep Learning}, Ch. 5.%
}

\newcommand{\AMLUnitVIISources}{%
Han, Kamber and Pei, \textit{Data Mining: Concepts and Techniques} (3e), Ch. 10.;%
James et al., \textit{An Introduction to Statistical Learning with Python}, Ch. 12.;%
M\"uller and Guido, \textit{Introduction to Machine Learning with Python}, \S3.5.;%
Ester, Kriegel, Sander and Xu, ``A density-based algorithm for discovering clusters (DBSCAN),'' KDD 1996.%
}

\newcommand{\AMLLabSources}{%
M\"uller and Guido, \textit{Introduction to Machine Learning with Python}, companion notebooks.;%
McKinney, \textit{Python for Data Analysis} (3e), O'Reilly, 2022.;%
Scikit-learn user guide, accessed 2026-08-04.;%
Matplotlib and Seaborn documentation, accessed 2026-08-04.%
}
%
  }{%
    % Question papers live at `private/<family>/<instrument>/`.
    \IfFileExists{../../../public/_shared/aml-sources.tex}{%
      % Source strings for Applied Machine Learning (CSAI2017P) — used by slides + notes.
% Prescribed texts and reference books are taken from the course syllabus.
% Support texts are the author-hosted open-access editions:
%   statlearning.com            (ISL with Python)
%   probml.github.io/pml-book   (Murphy, Probabilistic Machine Learning)
%   d2l.ai                      (Dive into Deep Learning)
%   mml-book.github.io          (Mathematics for Machine Learning)

\newcommand{\AMLRepoURL}{https://github.com/tali7c/Applied-Machine-Learning}

% --- prescribed -------------------------------------------------------------
\newcommand{\AMLPrescribed}{%
M\"uller and Guido, \textit{Introduction to Machine Learning with Python}, O'Reilly, 2016.;%
Bishop, \textit{Pattern Recognition and Machine Learning}, Springer, 2006.;%
Murphy, \textit{Machine Learning: A Probabilistic Perspective}, MIT Press, 2012.;%
Han, Kamber and Pei, \textit{Data Mining: Concepts and Techniques} (3e), Morgan Kaufmann, 2012.%
}

% --- unit-wise ---------------------------------------------------------------
\newcommand{\AMLUnitISources}{%
M\"uller and Guido, \textit{Introduction to Machine Learning with Python}, Ch. 1.;%
Bishop, \textit{Pattern Recognition and Machine Learning}, Ch. 1.;%
Murphy, \textit{Probabilistic Machine Learning: An Introduction}, MIT Press, 2022, Ch. 1.;%
Mitchell, \textit{Machine Learning}, McGraw Hill, 1997, Ch. 1.;%
Scikit-learn documentation, accessed 2026-08-04.%
}

\newcommand{\AMLUnitIISources}{%
Bishop, \textit{Pattern Recognition and Machine Learning}, \S1.5, \S4.3, \S7.1.;%
Zhang, Lipton, Li and Smola, \textit{Dive into Deep Learning}, \S3.4.;%
Shalev-Shwartz and Ben-David, \textit{Understanding Machine Learning}, Ch. 15.;%
Murphy, \textit{Probabilistic Machine Learning: An Introduction}, Ch. 5.%
}

\newcommand{\AMLUnitIIISources}{%
Zhang, Lipton, Li and Smola, \textit{Dive into Deep Learning}, Ch. 12 (Optimization Algorithms).;%
Deisenroth, Faisal and Ong, \textit{Mathematics for Machine Learning}, Ch. 7.;%
Kingma and Ba, ``Adam: A Method for Stochastic Optimization,'' ICLR 2015.;%
Ruder, ``An overview of gradient descent optimization algorithms,'' arXiv:1609.04747, 2016.%
}

\newcommand{\AMLUnitIVSources}{%
Han, Kamber and Pei, \textit{Data Mining: Concepts and Techniques} (3e), Ch. 3.;%
M\"uller and Guido, \textit{Introduction to Machine Learning with Python}, Ch. 3--4.;%
James, Witten, Hastie, Tibshirani and Taylor, \textit{An Introduction to Statistical Learning with Python}, Ch. 5--6, 12.;%
Scikit-learn documentation (preprocessing, model selection), accessed 2026-08-04.%
}

\newcommand{\AMLUnitVSources}{%
James et al., \textit{An Introduction to Statistical Learning with Python}, Ch. 3--6.;%
Bishop, \textit{Pattern Recognition and Machine Learning}, Ch. 3.;%
Deisenroth, Faisal and Ong, \textit{Mathematics for Machine Learning}, Ch. 9.;%
M\"uller and Guido, \textit{Introduction to Machine Learning with Python}, \S2.3.%
}

\newcommand{\AMLUnitVISources}{%
Han, Kamber and Pei, \textit{Data Mining: Concepts and Techniques} (3e), Ch. 8--9.;%
James et al., \textit{An Introduction to Statistical Learning with Python}, Ch. 8--9.;%
Bishop, \textit{Pattern Recognition and Machine Learning}, Ch. 4--5, 7, 14.;%
Zhang et al., \textit{Dive into Deep Learning}, Ch. 5.%
}

\newcommand{\AMLUnitVIISources}{%
Han, Kamber and Pei, \textit{Data Mining: Concepts and Techniques} (3e), Ch. 10.;%
James et al., \textit{An Introduction to Statistical Learning with Python}, Ch. 12.;%
M\"uller and Guido, \textit{Introduction to Machine Learning with Python}, \S3.5.;%
Ester, Kriegel, Sander and Xu, ``A density-based algorithm for discovering clusters (DBSCAN),'' KDD 1996.%
}

\newcommand{\AMLLabSources}{%
M\"uller and Guido, \textit{Introduction to Machine Learning with Python}, companion notebooks.;%
McKinney, \textit{Python for Data Analysis} (3e), O'Reilly, 2022.;%
Scikit-learn user guide, accessed 2026-08-04.;%
Matplotlib and Seaborn documentation, accessed 2026-08-04.%
}
%
    }{%
      % Marking schemes live at `private/solutions/<family>/<id>/latex/`.
      \IfFileExists{../../../../../public/_shared/aml-sources.tex}{%
        % Source strings for Applied Machine Learning (CSAI2017P) — used by slides + notes.
% Prescribed texts and reference books are taken from the course syllabus.
% Support texts are the author-hosted open-access editions:
%   statlearning.com            (ISL with Python)
%   probml.github.io/pml-book   (Murphy, Probabilistic Machine Learning)
%   d2l.ai                      (Dive into Deep Learning)
%   mml-book.github.io          (Mathematics for Machine Learning)

\newcommand{\AMLRepoURL}{https://github.com/tali7c/Applied-Machine-Learning}

% --- prescribed -------------------------------------------------------------
\newcommand{\AMLPrescribed}{%
M\"uller and Guido, \textit{Introduction to Machine Learning with Python}, O'Reilly, 2016.;%
Bishop, \textit{Pattern Recognition and Machine Learning}, Springer, 2006.;%
Murphy, \textit{Machine Learning: A Probabilistic Perspective}, MIT Press, 2012.;%
Han, Kamber and Pei, \textit{Data Mining: Concepts and Techniques} (3e), Morgan Kaufmann, 2012.%
}

% --- unit-wise ---------------------------------------------------------------
\newcommand{\AMLUnitISources}{%
M\"uller and Guido, \textit{Introduction to Machine Learning with Python}, Ch. 1.;%
Bishop, \textit{Pattern Recognition and Machine Learning}, Ch. 1.;%
Murphy, \textit{Probabilistic Machine Learning: An Introduction}, MIT Press, 2022, Ch. 1.;%
Mitchell, \textit{Machine Learning}, McGraw Hill, 1997, Ch. 1.;%
Scikit-learn documentation, accessed 2026-08-04.%
}

\newcommand{\AMLUnitIISources}{%
Bishop, \textit{Pattern Recognition and Machine Learning}, \S1.5, \S4.3, \S7.1.;%
Zhang, Lipton, Li and Smola, \textit{Dive into Deep Learning}, \S3.4.;%
Shalev-Shwartz and Ben-David, \textit{Understanding Machine Learning}, Ch. 15.;%
Murphy, \textit{Probabilistic Machine Learning: An Introduction}, Ch. 5.%
}

\newcommand{\AMLUnitIIISources}{%
Zhang, Lipton, Li and Smola, \textit{Dive into Deep Learning}, Ch. 12 (Optimization Algorithms).;%
Deisenroth, Faisal and Ong, \textit{Mathematics for Machine Learning}, Ch. 7.;%
Kingma and Ba, ``Adam: A Method for Stochastic Optimization,'' ICLR 2015.;%
Ruder, ``An overview of gradient descent optimization algorithms,'' arXiv:1609.04747, 2016.%
}

\newcommand{\AMLUnitIVSources}{%
Han, Kamber and Pei, \textit{Data Mining: Concepts and Techniques} (3e), Ch. 3.;%
M\"uller and Guido, \textit{Introduction to Machine Learning with Python}, Ch. 3--4.;%
James, Witten, Hastie, Tibshirani and Taylor, \textit{An Introduction to Statistical Learning with Python}, Ch. 5--6, 12.;%
Scikit-learn documentation (preprocessing, model selection), accessed 2026-08-04.%
}

\newcommand{\AMLUnitVSources}{%
James et al., \textit{An Introduction to Statistical Learning with Python}, Ch. 3--6.;%
Bishop, \textit{Pattern Recognition and Machine Learning}, Ch. 3.;%
Deisenroth, Faisal and Ong, \textit{Mathematics for Machine Learning}, Ch. 9.;%
M\"uller and Guido, \textit{Introduction to Machine Learning with Python}, \S2.3.%
}

\newcommand{\AMLUnitVISources}{%
Han, Kamber and Pei, \textit{Data Mining: Concepts and Techniques} (3e), Ch. 8--9.;%
James et al., \textit{An Introduction to Statistical Learning with Python}, Ch. 8--9.;%
Bishop, \textit{Pattern Recognition and Machine Learning}, Ch. 4--5, 7, 14.;%
Zhang et al., \textit{Dive into Deep Learning}, Ch. 5.%
}

\newcommand{\AMLUnitVIISources}{%
Han, Kamber and Pei, \textit{Data Mining: Concepts and Techniques} (3e), Ch. 10.;%
James et al., \textit{An Introduction to Statistical Learning with Python}, Ch. 12.;%
M\"uller and Guido, \textit{Introduction to Machine Learning with Python}, \S3.5.;%
Ester, Kriegel, Sander and Xu, ``A density-based algorithm for discovering clusters (DBSCAN),'' KDD 1996.%
}

\newcommand{\AMLLabSources}{%
M\"uller and Guido, \textit{Introduction to Machine Learning with Python}, companion notebooks.;%
McKinney, \textit{Python for Data Analysis} (3e), O'Reilly, 2022.;%
Scikit-learn user guide, accessed 2026-08-04.;%
Matplotlib and Seaborn documentation, accessed 2026-08-04.%
}
%
      }{%
        % Source strings for Applied Machine Learning (CSAI2017P) — used by slides + notes.
% Prescribed texts and reference books are taken from the course syllabus.
% Support texts are the author-hosted open-access editions:
%   statlearning.com            (ISL with Python)
%   probml.github.io/pml-book   (Murphy, Probabilistic Machine Learning)
%   d2l.ai                      (Dive into Deep Learning)
%   mml-book.github.io          (Mathematics for Machine Learning)

\newcommand{\AMLRepoURL}{https://github.com/tali7c/Applied-Machine-Learning}

% --- prescribed -------------------------------------------------------------
\newcommand{\AMLPrescribed}{%
M\"uller and Guido, \textit{Introduction to Machine Learning with Python}, O'Reilly, 2016.;%
Bishop, \textit{Pattern Recognition and Machine Learning}, Springer, 2006.;%
Murphy, \textit{Machine Learning: A Probabilistic Perspective}, MIT Press, 2012.;%
Han, Kamber and Pei, \textit{Data Mining: Concepts and Techniques} (3e), Morgan Kaufmann, 2012.%
}

% --- unit-wise ---------------------------------------------------------------
\newcommand{\AMLUnitISources}{%
M\"uller and Guido, \textit{Introduction to Machine Learning with Python}, Ch. 1.;%
Bishop, \textit{Pattern Recognition and Machine Learning}, Ch. 1.;%
Murphy, \textit{Probabilistic Machine Learning: An Introduction}, MIT Press, 2022, Ch. 1.;%
Mitchell, \textit{Machine Learning}, McGraw Hill, 1997, Ch. 1.;%
Scikit-learn documentation, accessed 2026-08-04.%
}

\newcommand{\AMLUnitIISources}{%
Bishop, \textit{Pattern Recognition and Machine Learning}, \S1.5, \S4.3, \S7.1.;%
Zhang, Lipton, Li and Smola, \textit{Dive into Deep Learning}, \S3.4.;%
Shalev-Shwartz and Ben-David, \textit{Understanding Machine Learning}, Ch. 15.;%
Murphy, \textit{Probabilistic Machine Learning: An Introduction}, Ch. 5.%
}

\newcommand{\AMLUnitIIISources}{%
Zhang, Lipton, Li and Smola, \textit{Dive into Deep Learning}, Ch. 12 (Optimization Algorithms).;%
Deisenroth, Faisal and Ong, \textit{Mathematics for Machine Learning}, Ch. 7.;%
Kingma and Ba, ``Adam: A Method for Stochastic Optimization,'' ICLR 2015.;%
Ruder, ``An overview of gradient descent optimization algorithms,'' arXiv:1609.04747, 2016.%
}

\newcommand{\AMLUnitIVSources}{%
Han, Kamber and Pei, \textit{Data Mining: Concepts and Techniques} (3e), Ch. 3.;%
M\"uller and Guido, \textit{Introduction to Machine Learning with Python}, Ch. 3--4.;%
James, Witten, Hastie, Tibshirani and Taylor, \textit{An Introduction to Statistical Learning with Python}, Ch. 5--6, 12.;%
Scikit-learn documentation (preprocessing, model selection), accessed 2026-08-04.%
}

\newcommand{\AMLUnitVSources}{%
James et al., \textit{An Introduction to Statistical Learning with Python}, Ch. 3--6.;%
Bishop, \textit{Pattern Recognition and Machine Learning}, Ch. 3.;%
Deisenroth, Faisal and Ong, \textit{Mathematics for Machine Learning}, Ch. 9.;%
M\"uller and Guido, \textit{Introduction to Machine Learning with Python}, \S2.3.%
}

\newcommand{\AMLUnitVISources}{%
Han, Kamber and Pei, \textit{Data Mining: Concepts and Techniques} (3e), Ch. 8--9.;%
James et al., \textit{An Introduction to Statistical Learning with Python}, Ch. 8--9.;%
Bishop, \textit{Pattern Recognition and Machine Learning}, Ch. 4--5, 7, 14.;%
Zhang et al., \textit{Dive into Deep Learning}, Ch. 5.%
}

\newcommand{\AMLUnitVIISources}{%
Han, Kamber and Pei, \textit{Data Mining: Concepts and Techniques} (3e), Ch. 10.;%
James et al., \textit{An Introduction to Statistical Learning with Python}, Ch. 12.;%
M\"uller and Guido, \textit{Introduction to Machine Learning with Python}, \S3.5.;%
Ester, Kriegel, Sander and Xu, ``A density-based algorithm for discovering clusters (DBSCAN),'' KDD 1996.%
}

\newcommand{\AMLLabSources}{%
M\"uller and Guido, \textit{Introduction to Machine Learning with Python}, companion notebooks.;%
McKinney, \textit{Python for Data Analysis} (3e), O'Reilly, 2022.;%
Scikit-learn user guide, accessed 2026-08-04.;%
Matplotlib and Seaborn documentation, accessed 2026-08-04.%
}
%
      }%
    }%
  }%
}
%
  }{%
    % Shared notes settings for Applied Machine Learning lectures (UPES).
% Keep this file free of \documentclass to allow reuse via \input.

\usepackage[utf8]{inputenc}
\usepackage[T1]{fontenc}
\usepackage{lmodern}          % scalable T1 fonts; without this pdflatex falls back
\input{glyphtounicode}        % to bitmapped EC Type 3 and the PDFs are neither
\pdfgentounicode=1            % searchable nor screen-reader friendly
\usepackage[a4paper,margin=1in]{geometry}
\usepackage{hyperref}
\usepackage{xurl}
\usepackage{booktabs}
\usepackage{amsmath}
\usepackage{amssymb}
\usepackage{listings}
\usepackage{xcolor}
\usepackage{enumitem}
\usepackage{parskip}

% Code listings (Python-oriented for this subject).
\lstset{
  language=Python,
  basicstyle=\ttfamily\small,
  keywordstyle=\color{blue},
  commentstyle=\color{gray},
  stringstyle=\color{teal},
  showstringspaces=false,
  columns=fullflexible,
  frame=single,
  framerule=0.2pt,
  breaklines=true
}

\setlist[itemize]{topsep=0.3em,itemsep=0.2em}
\setlist[enumerate]{topsep=0.3em,itemsep=0.2em}

% Simple per-section source line.
\newcommand{\notesources}[1]{\par\noindent{\small\color{gray}\textit{Sources:} \sloppy #1\par}}

% Unit-wise source strings for this subject (used by slides + notes).
% Path is relative to the lecture `latex/` folder (the compilation CWD).
\IfFileExists{../../../_shared/aml-sources.tex}{%
  % Source strings for Applied Machine Learning (CSAI2017P) — used by slides + notes.
% Prescribed texts and reference books are taken from the course syllabus.
% Support texts are the author-hosted open-access editions:
%   statlearning.com            (ISL with Python)
%   probml.github.io/pml-book   (Murphy, Probabilistic Machine Learning)
%   d2l.ai                      (Dive into Deep Learning)
%   mml-book.github.io          (Mathematics for Machine Learning)

\newcommand{\AMLRepoURL}{https://github.com/tali7c/Applied-Machine-Learning}

% --- prescribed -------------------------------------------------------------
\newcommand{\AMLPrescribed}{%
M\"uller and Guido, \textit{Introduction to Machine Learning with Python}, O'Reilly, 2016.;%
Bishop, \textit{Pattern Recognition and Machine Learning}, Springer, 2006.;%
Murphy, \textit{Machine Learning: A Probabilistic Perspective}, MIT Press, 2012.;%
Han, Kamber and Pei, \textit{Data Mining: Concepts and Techniques} (3e), Morgan Kaufmann, 2012.%
}

% --- unit-wise ---------------------------------------------------------------
\newcommand{\AMLUnitISources}{%
M\"uller and Guido, \textit{Introduction to Machine Learning with Python}, Ch. 1.;%
Bishop, \textit{Pattern Recognition and Machine Learning}, Ch. 1.;%
Murphy, \textit{Probabilistic Machine Learning: An Introduction}, MIT Press, 2022, Ch. 1.;%
Mitchell, \textit{Machine Learning}, McGraw Hill, 1997, Ch. 1.;%
Scikit-learn documentation, accessed 2026-08-04.%
}

\newcommand{\AMLUnitIISources}{%
Bishop, \textit{Pattern Recognition and Machine Learning}, \S1.5, \S4.3, \S7.1.;%
Zhang, Lipton, Li and Smola, \textit{Dive into Deep Learning}, \S3.4.;%
Shalev-Shwartz and Ben-David, \textit{Understanding Machine Learning}, Ch. 15.;%
Murphy, \textit{Probabilistic Machine Learning: An Introduction}, Ch. 5.%
}

\newcommand{\AMLUnitIIISources}{%
Zhang, Lipton, Li and Smola, \textit{Dive into Deep Learning}, Ch. 12 (Optimization Algorithms).;%
Deisenroth, Faisal and Ong, \textit{Mathematics for Machine Learning}, Ch. 7.;%
Kingma and Ba, ``Adam: A Method for Stochastic Optimization,'' ICLR 2015.;%
Ruder, ``An overview of gradient descent optimization algorithms,'' arXiv:1609.04747, 2016.%
}

\newcommand{\AMLUnitIVSources}{%
Han, Kamber and Pei, \textit{Data Mining: Concepts and Techniques} (3e), Ch. 3.;%
M\"uller and Guido, \textit{Introduction to Machine Learning with Python}, Ch. 3--4.;%
James, Witten, Hastie, Tibshirani and Taylor, \textit{An Introduction to Statistical Learning with Python}, Ch. 5--6, 12.;%
Scikit-learn documentation (preprocessing, model selection), accessed 2026-08-04.%
}

\newcommand{\AMLUnitVSources}{%
James et al., \textit{An Introduction to Statistical Learning with Python}, Ch. 3--6.;%
Bishop, \textit{Pattern Recognition and Machine Learning}, Ch. 3.;%
Deisenroth, Faisal and Ong, \textit{Mathematics for Machine Learning}, Ch. 9.;%
M\"uller and Guido, \textit{Introduction to Machine Learning with Python}, \S2.3.%
}

\newcommand{\AMLUnitVISources}{%
Han, Kamber and Pei, \textit{Data Mining: Concepts and Techniques} (3e), Ch. 8--9.;%
James et al., \textit{An Introduction to Statistical Learning with Python}, Ch. 8--9.;%
Bishop, \textit{Pattern Recognition and Machine Learning}, Ch. 4--5, 7, 14.;%
Zhang et al., \textit{Dive into Deep Learning}, Ch. 5.%
}

\newcommand{\AMLUnitVIISources}{%
Han, Kamber and Pei, \textit{Data Mining: Concepts and Techniques} (3e), Ch. 10.;%
James et al., \textit{An Introduction to Statistical Learning with Python}, Ch. 12.;%
M\"uller and Guido, \textit{Introduction to Machine Learning with Python}, \S3.5.;%
Ester, Kriegel, Sander and Xu, ``A density-based algorithm for discovering clusters (DBSCAN),'' KDD 1996.%
}

\newcommand{\AMLLabSources}{%
M\"uller and Guido, \textit{Introduction to Machine Learning with Python}, companion notebooks.;%
McKinney, \textit{Python for Data Analysis} (3e), O'Reilly, 2022.;%
Scikit-learn user guide, accessed 2026-08-04.;%
Matplotlib and Seaborn documentation, accessed 2026-08-04.%
}
%
}{%
  % Faculty copies live under `private/` so their compile CWD is different.
  \IfFileExists{../../../../public/_shared/aml-sources.tex}{%
    % Source strings for Applied Machine Learning (CSAI2017P) — used by slides + notes.
% Prescribed texts and reference books are taken from the course syllabus.
% Support texts are the author-hosted open-access editions:
%   statlearning.com            (ISL with Python)
%   probml.github.io/pml-book   (Murphy, Probabilistic Machine Learning)
%   d2l.ai                      (Dive into Deep Learning)
%   mml-book.github.io          (Mathematics for Machine Learning)

\newcommand{\AMLRepoURL}{https://github.com/tali7c/Applied-Machine-Learning}

% --- prescribed -------------------------------------------------------------
\newcommand{\AMLPrescribed}{%
M\"uller and Guido, \textit{Introduction to Machine Learning with Python}, O'Reilly, 2016.;%
Bishop, \textit{Pattern Recognition and Machine Learning}, Springer, 2006.;%
Murphy, \textit{Machine Learning: A Probabilistic Perspective}, MIT Press, 2012.;%
Han, Kamber and Pei, \textit{Data Mining: Concepts and Techniques} (3e), Morgan Kaufmann, 2012.%
}

% --- unit-wise ---------------------------------------------------------------
\newcommand{\AMLUnitISources}{%
M\"uller and Guido, \textit{Introduction to Machine Learning with Python}, Ch. 1.;%
Bishop, \textit{Pattern Recognition and Machine Learning}, Ch. 1.;%
Murphy, \textit{Probabilistic Machine Learning: An Introduction}, MIT Press, 2022, Ch. 1.;%
Mitchell, \textit{Machine Learning}, McGraw Hill, 1997, Ch. 1.;%
Scikit-learn documentation, accessed 2026-08-04.%
}

\newcommand{\AMLUnitIISources}{%
Bishop, \textit{Pattern Recognition and Machine Learning}, \S1.5, \S4.3, \S7.1.;%
Zhang, Lipton, Li and Smola, \textit{Dive into Deep Learning}, \S3.4.;%
Shalev-Shwartz and Ben-David, \textit{Understanding Machine Learning}, Ch. 15.;%
Murphy, \textit{Probabilistic Machine Learning: An Introduction}, Ch. 5.%
}

\newcommand{\AMLUnitIIISources}{%
Zhang, Lipton, Li and Smola, \textit{Dive into Deep Learning}, Ch. 12 (Optimization Algorithms).;%
Deisenroth, Faisal and Ong, \textit{Mathematics for Machine Learning}, Ch. 7.;%
Kingma and Ba, ``Adam: A Method for Stochastic Optimization,'' ICLR 2015.;%
Ruder, ``An overview of gradient descent optimization algorithms,'' arXiv:1609.04747, 2016.%
}

\newcommand{\AMLUnitIVSources}{%
Han, Kamber and Pei, \textit{Data Mining: Concepts and Techniques} (3e), Ch. 3.;%
M\"uller and Guido, \textit{Introduction to Machine Learning with Python}, Ch. 3--4.;%
James, Witten, Hastie, Tibshirani and Taylor, \textit{An Introduction to Statistical Learning with Python}, Ch. 5--6, 12.;%
Scikit-learn documentation (preprocessing, model selection), accessed 2026-08-04.%
}

\newcommand{\AMLUnitVSources}{%
James et al., \textit{An Introduction to Statistical Learning with Python}, Ch. 3--6.;%
Bishop, \textit{Pattern Recognition and Machine Learning}, Ch. 3.;%
Deisenroth, Faisal and Ong, \textit{Mathematics for Machine Learning}, Ch. 9.;%
M\"uller and Guido, \textit{Introduction to Machine Learning with Python}, \S2.3.%
}

\newcommand{\AMLUnitVISources}{%
Han, Kamber and Pei, \textit{Data Mining: Concepts and Techniques} (3e), Ch. 8--9.;%
James et al., \textit{An Introduction to Statistical Learning with Python}, Ch. 8--9.;%
Bishop, \textit{Pattern Recognition and Machine Learning}, Ch. 4--5, 7, 14.;%
Zhang et al., \textit{Dive into Deep Learning}, Ch. 5.%
}

\newcommand{\AMLUnitVIISources}{%
Han, Kamber and Pei, \textit{Data Mining: Concepts and Techniques} (3e), Ch. 10.;%
James et al., \textit{An Introduction to Statistical Learning with Python}, Ch. 12.;%
M\"uller and Guido, \textit{Introduction to Machine Learning with Python}, \S3.5.;%
Ester, Kriegel, Sander and Xu, ``A density-based algorithm for discovering clusters (DBSCAN),'' KDD 1996.%
}

\newcommand{\AMLLabSources}{%
M\"uller and Guido, \textit{Introduction to Machine Learning with Python}, companion notebooks.;%
McKinney, \textit{Python for Data Analysis} (3e), O'Reilly, 2022.;%
Scikit-learn user guide, accessed 2026-08-04.;%
Matplotlib and Seaborn documentation, accessed 2026-08-04.%
}
%
  }{%
    % Question papers live at `private/<family>/<instrument>/`.
    \IfFileExists{../../../public/_shared/aml-sources.tex}{%
      % Source strings for Applied Machine Learning (CSAI2017P) — used by slides + notes.
% Prescribed texts and reference books are taken from the course syllabus.
% Support texts are the author-hosted open-access editions:
%   statlearning.com            (ISL with Python)
%   probml.github.io/pml-book   (Murphy, Probabilistic Machine Learning)
%   d2l.ai                      (Dive into Deep Learning)
%   mml-book.github.io          (Mathematics for Machine Learning)

\newcommand{\AMLRepoURL}{https://github.com/tali7c/Applied-Machine-Learning}

% --- prescribed -------------------------------------------------------------
\newcommand{\AMLPrescribed}{%
M\"uller and Guido, \textit{Introduction to Machine Learning with Python}, O'Reilly, 2016.;%
Bishop, \textit{Pattern Recognition and Machine Learning}, Springer, 2006.;%
Murphy, \textit{Machine Learning: A Probabilistic Perspective}, MIT Press, 2012.;%
Han, Kamber and Pei, \textit{Data Mining: Concepts and Techniques} (3e), Morgan Kaufmann, 2012.%
}

% --- unit-wise ---------------------------------------------------------------
\newcommand{\AMLUnitISources}{%
M\"uller and Guido, \textit{Introduction to Machine Learning with Python}, Ch. 1.;%
Bishop, \textit{Pattern Recognition and Machine Learning}, Ch. 1.;%
Murphy, \textit{Probabilistic Machine Learning: An Introduction}, MIT Press, 2022, Ch. 1.;%
Mitchell, \textit{Machine Learning}, McGraw Hill, 1997, Ch. 1.;%
Scikit-learn documentation, accessed 2026-08-04.%
}

\newcommand{\AMLUnitIISources}{%
Bishop, \textit{Pattern Recognition and Machine Learning}, \S1.5, \S4.3, \S7.1.;%
Zhang, Lipton, Li and Smola, \textit{Dive into Deep Learning}, \S3.4.;%
Shalev-Shwartz and Ben-David, \textit{Understanding Machine Learning}, Ch. 15.;%
Murphy, \textit{Probabilistic Machine Learning: An Introduction}, Ch. 5.%
}

\newcommand{\AMLUnitIIISources}{%
Zhang, Lipton, Li and Smola, \textit{Dive into Deep Learning}, Ch. 12 (Optimization Algorithms).;%
Deisenroth, Faisal and Ong, \textit{Mathematics for Machine Learning}, Ch. 7.;%
Kingma and Ba, ``Adam: A Method for Stochastic Optimization,'' ICLR 2015.;%
Ruder, ``An overview of gradient descent optimization algorithms,'' arXiv:1609.04747, 2016.%
}

\newcommand{\AMLUnitIVSources}{%
Han, Kamber and Pei, \textit{Data Mining: Concepts and Techniques} (3e), Ch. 3.;%
M\"uller and Guido, \textit{Introduction to Machine Learning with Python}, Ch. 3--4.;%
James, Witten, Hastie, Tibshirani and Taylor, \textit{An Introduction to Statistical Learning with Python}, Ch. 5--6, 12.;%
Scikit-learn documentation (preprocessing, model selection), accessed 2026-08-04.%
}

\newcommand{\AMLUnitVSources}{%
James et al., \textit{An Introduction to Statistical Learning with Python}, Ch. 3--6.;%
Bishop, \textit{Pattern Recognition and Machine Learning}, Ch. 3.;%
Deisenroth, Faisal and Ong, \textit{Mathematics for Machine Learning}, Ch. 9.;%
M\"uller and Guido, \textit{Introduction to Machine Learning with Python}, \S2.3.%
}

\newcommand{\AMLUnitVISources}{%
Han, Kamber and Pei, \textit{Data Mining: Concepts and Techniques} (3e), Ch. 8--9.;%
James et al., \textit{An Introduction to Statistical Learning with Python}, Ch. 8--9.;%
Bishop, \textit{Pattern Recognition and Machine Learning}, Ch. 4--5, 7, 14.;%
Zhang et al., \textit{Dive into Deep Learning}, Ch. 5.%
}

\newcommand{\AMLUnitVIISources}{%
Han, Kamber and Pei, \textit{Data Mining: Concepts and Techniques} (3e), Ch. 10.;%
James et al., \textit{An Introduction to Statistical Learning with Python}, Ch. 12.;%
M\"uller and Guido, \textit{Introduction to Machine Learning with Python}, \S3.5.;%
Ester, Kriegel, Sander and Xu, ``A density-based algorithm for discovering clusters (DBSCAN),'' KDD 1996.%
}

\newcommand{\AMLLabSources}{%
M\"uller and Guido, \textit{Introduction to Machine Learning with Python}, companion notebooks.;%
McKinney, \textit{Python for Data Analysis} (3e), O'Reilly, 2022.;%
Scikit-learn user guide, accessed 2026-08-04.;%
Matplotlib and Seaborn documentation, accessed 2026-08-04.%
}
%
    }{%
      % Marking schemes live at `private/solutions/<family>/<id>/latex/`.
      \IfFileExists{../../../../../public/_shared/aml-sources.tex}{%
        % Source strings for Applied Machine Learning (CSAI2017P) — used by slides + notes.
% Prescribed texts and reference books are taken from the course syllabus.
% Support texts are the author-hosted open-access editions:
%   statlearning.com            (ISL with Python)
%   probml.github.io/pml-book   (Murphy, Probabilistic Machine Learning)
%   d2l.ai                      (Dive into Deep Learning)
%   mml-book.github.io          (Mathematics for Machine Learning)

\newcommand{\AMLRepoURL}{https://github.com/tali7c/Applied-Machine-Learning}

% --- prescribed -------------------------------------------------------------
\newcommand{\AMLPrescribed}{%
M\"uller and Guido, \textit{Introduction to Machine Learning with Python}, O'Reilly, 2016.;%
Bishop, \textit{Pattern Recognition and Machine Learning}, Springer, 2006.;%
Murphy, \textit{Machine Learning: A Probabilistic Perspective}, MIT Press, 2012.;%
Han, Kamber and Pei, \textit{Data Mining: Concepts and Techniques} (3e), Morgan Kaufmann, 2012.%
}

% --- unit-wise ---------------------------------------------------------------
\newcommand{\AMLUnitISources}{%
M\"uller and Guido, \textit{Introduction to Machine Learning with Python}, Ch. 1.;%
Bishop, \textit{Pattern Recognition and Machine Learning}, Ch. 1.;%
Murphy, \textit{Probabilistic Machine Learning: An Introduction}, MIT Press, 2022, Ch. 1.;%
Mitchell, \textit{Machine Learning}, McGraw Hill, 1997, Ch. 1.;%
Scikit-learn documentation, accessed 2026-08-04.%
}

\newcommand{\AMLUnitIISources}{%
Bishop, \textit{Pattern Recognition and Machine Learning}, \S1.5, \S4.3, \S7.1.;%
Zhang, Lipton, Li and Smola, \textit{Dive into Deep Learning}, \S3.4.;%
Shalev-Shwartz and Ben-David, \textit{Understanding Machine Learning}, Ch. 15.;%
Murphy, \textit{Probabilistic Machine Learning: An Introduction}, Ch. 5.%
}

\newcommand{\AMLUnitIIISources}{%
Zhang, Lipton, Li and Smola, \textit{Dive into Deep Learning}, Ch. 12 (Optimization Algorithms).;%
Deisenroth, Faisal and Ong, \textit{Mathematics for Machine Learning}, Ch. 7.;%
Kingma and Ba, ``Adam: A Method for Stochastic Optimization,'' ICLR 2015.;%
Ruder, ``An overview of gradient descent optimization algorithms,'' arXiv:1609.04747, 2016.%
}

\newcommand{\AMLUnitIVSources}{%
Han, Kamber and Pei, \textit{Data Mining: Concepts and Techniques} (3e), Ch. 3.;%
M\"uller and Guido, \textit{Introduction to Machine Learning with Python}, Ch. 3--4.;%
James, Witten, Hastie, Tibshirani and Taylor, \textit{An Introduction to Statistical Learning with Python}, Ch. 5--6, 12.;%
Scikit-learn documentation (preprocessing, model selection), accessed 2026-08-04.%
}

\newcommand{\AMLUnitVSources}{%
James et al., \textit{An Introduction to Statistical Learning with Python}, Ch. 3--6.;%
Bishop, \textit{Pattern Recognition and Machine Learning}, Ch. 3.;%
Deisenroth, Faisal and Ong, \textit{Mathematics for Machine Learning}, Ch. 9.;%
M\"uller and Guido, \textit{Introduction to Machine Learning with Python}, \S2.3.%
}

\newcommand{\AMLUnitVISources}{%
Han, Kamber and Pei, \textit{Data Mining: Concepts and Techniques} (3e), Ch. 8--9.;%
James et al., \textit{An Introduction to Statistical Learning with Python}, Ch. 8--9.;%
Bishop, \textit{Pattern Recognition and Machine Learning}, Ch. 4--5, 7, 14.;%
Zhang et al., \textit{Dive into Deep Learning}, Ch. 5.%
}

\newcommand{\AMLUnitVIISources}{%
Han, Kamber and Pei, \textit{Data Mining: Concepts and Techniques} (3e), Ch. 10.;%
James et al., \textit{An Introduction to Statistical Learning with Python}, Ch. 12.;%
M\"uller and Guido, \textit{Introduction to Machine Learning with Python}, \S3.5.;%
Ester, Kriegel, Sander and Xu, ``A density-based algorithm for discovering clusters (DBSCAN),'' KDD 1996.%
}

\newcommand{\AMLLabSources}{%
M\"uller and Guido, \textit{Introduction to Machine Learning with Python}, companion notebooks.;%
McKinney, \textit{Python for Data Analysis} (3e), O'Reilly, 2022.;%
Scikit-learn user guide, accessed 2026-08-04.;%
Matplotlib and Seaborn documentation, accessed 2026-08-04.%
}
%
      }{%
        % Source strings for Applied Machine Learning (CSAI2017P) — used by slides + notes.
% Prescribed texts and reference books are taken from the course syllabus.
% Support texts are the author-hosted open-access editions:
%   statlearning.com            (ISL with Python)
%   probml.github.io/pml-book   (Murphy, Probabilistic Machine Learning)
%   d2l.ai                      (Dive into Deep Learning)
%   mml-book.github.io          (Mathematics for Machine Learning)

\newcommand{\AMLRepoURL}{https://github.com/tali7c/Applied-Machine-Learning}

% --- prescribed -------------------------------------------------------------
\newcommand{\AMLPrescribed}{%
M\"uller and Guido, \textit{Introduction to Machine Learning with Python}, O'Reilly, 2016.;%
Bishop, \textit{Pattern Recognition and Machine Learning}, Springer, 2006.;%
Murphy, \textit{Machine Learning: A Probabilistic Perspective}, MIT Press, 2012.;%
Han, Kamber and Pei, \textit{Data Mining: Concepts and Techniques} (3e), Morgan Kaufmann, 2012.%
}

% --- unit-wise ---------------------------------------------------------------
\newcommand{\AMLUnitISources}{%
M\"uller and Guido, \textit{Introduction to Machine Learning with Python}, Ch. 1.;%
Bishop, \textit{Pattern Recognition and Machine Learning}, Ch. 1.;%
Murphy, \textit{Probabilistic Machine Learning: An Introduction}, MIT Press, 2022, Ch. 1.;%
Mitchell, \textit{Machine Learning}, McGraw Hill, 1997, Ch. 1.;%
Scikit-learn documentation, accessed 2026-08-04.%
}

\newcommand{\AMLUnitIISources}{%
Bishop, \textit{Pattern Recognition and Machine Learning}, \S1.5, \S4.3, \S7.1.;%
Zhang, Lipton, Li and Smola, \textit{Dive into Deep Learning}, \S3.4.;%
Shalev-Shwartz and Ben-David, \textit{Understanding Machine Learning}, Ch. 15.;%
Murphy, \textit{Probabilistic Machine Learning: An Introduction}, Ch. 5.%
}

\newcommand{\AMLUnitIIISources}{%
Zhang, Lipton, Li and Smola, \textit{Dive into Deep Learning}, Ch. 12 (Optimization Algorithms).;%
Deisenroth, Faisal and Ong, \textit{Mathematics for Machine Learning}, Ch. 7.;%
Kingma and Ba, ``Adam: A Method for Stochastic Optimization,'' ICLR 2015.;%
Ruder, ``An overview of gradient descent optimization algorithms,'' arXiv:1609.04747, 2016.%
}

\newcommand{\AMLUnitIVSources}{%
Han, Kamber and Pei, \textit{Data Mining: Concepts and Techniques} (3e), Ch. 3.;%
M\"uller and Guido, \textit{Introduction to Machine Learning with Python}, Ch. 3--4.;%
James, Witten, Hastie, Tibshirani and Taylor, \textit{An Introduction to Statistical Learning with Python}, Ch. 5--6, 12.;%
Scikit-learn documentation (preprocessing, model selection), accessed 2026-08-04.%
}

\newcommand{\AMLUnitVSources}{%
James et al., \textit{An Introduction to Statistical Learning with Python}, Ch. 3--6.;%
Bishop, \textit{Pattern Recognition and Machine Learning}, Ch. 3.;%
Deisenroth, Faisal and Ong, \textit{Mathematics for Machine Learning}, Ch. 9.;%
M\"uller and Guido, \textit{Introduction to Machine Learning with Python}, \S2.3.%
}

\newcommand{\AMLUnitVISources}{%
Han, Kamber and Pei, \textit{Data Mining: Concepts and Techniques} (3e), Ch. 8--9.;%
James et al., \textit{An Introduction to Statistical Learning with Python}, Ch. 8--9.;%
Bishop, \textit{Pattern Recognition and Machine Learning}, Ch. 4--5, 7, 14.;%
Zhang et al., \textit{Dive into Deep Learning}, Ch. 5.%
}

\newcommand{\AMLUnitVIISources}{%
Han, Kamber and Pei, \textit{Data Mining: Concepts and Techniques} (3e), Ch. 10.;%
James et al., \textit{An Introduction to Statistical Learning with Python}, Ch. 12.;%
M\"uller and Guido, \textit{Introduction to Machine Learning with Python}, \S3.5.;%
Ester, Kriegel, Sander and Xu, ``A density-based algorithm for discovering clusters (DBSCAN),'' KDD 1996.%
}

\newcommand{\AMLLabSources}{%
M\"uller and Guido, \textit{Introduction to Machine Learning with Python}, companion notebooks.;%
McKinney, \textit{Python for Data Analysis} (3e), O'Reilly, 2022.;%
Scikit-learn user guide, accessed 2026-08-04.;%
Matplotlib and Seaborn documentation, accessed 2026-08-04.%
}
%
      }%
    }%
  }%
}
%
  }%
}

\usepackage{fancyhdr}
\usepackage{tcolorbox}
\tcbuselibrary{skins,breakable}

\usepackage{array}
\usepackage[htt]{hyphenat}   % allow \texttt{...} to hyphenate, else long
                             % identifiers overflow the measure

\hypersetup{colorlinks=true, linkcolor=blue, urlcolor=blue, breaklinks=true}
\setlist{itemsep=0.35em, topsep=0.35em}
\sloppy
\emergencystretch=2em

\providecommand{\CourseCode}{CSAI2017P: Applied Machine Learning (Lab)}
\providecommand{\AssignmentName}{Lab Assignment}

\pagestyle{fancy}
\fancyhf{}
\lhead{\small\CourseCode}
\rhead{\small\AssignmentName}
\cfoot{\thepage}
\setlength{\headheight}{14pt}

% Per-question mark tag, right aligned.
\newcommand{\QMarks}[1]{\hfill \textbf{[#1 marks]}}

% Title block for a brief.
\newcommand{\LAtitle}[4]{%
  \begin{center}
    {\LARGE \textbf{#1}}\\[0.25em]
    {\large \CourseCode}\\[0.25em]
    \normalsize #2
  \end{center}
  \noindent
  \textbf{Instructor:} Dr.\ Tofik Ali \hfill \textbf{Due:} #3\\
  \textbf{Student Name:} \rule{5cm}{0.4pt} \hfill \textbf{SAP ID:} \rule{3.5cm}{0.4pt}\\
  \textbf{Batch:} \rule{2.5cm}{0.4pt} \hfill \textbf{Lab quiz:} #4
  \vspace{0.6em}\hrule\vspace{0.8em}}

% Title block for a marking scheme (no student fields).
\newcommand{\MStitle}[3]{%
  \begin{center}
    {\LARGE \textbf{#1}}\\[0.25em]
    {\large \CourseCode}\\[0.25em]
    \normalsize #2\\[0.4em]
    {\color{red}\textbf{MARKING SCHEME --- NOT FOR CIRCULATION}}
  \end{center}
  \noindent \textbf{Instructor:} Dr.\ Tofik Ali \hfill \textbf{Assignment due:} #3
  \vspace{0.6em}\hrule\vspace{0.8em}}

\newtcolorbox{infobox}[1][Note]{
  breakable, enhanced, colback=gray!6, colframe=gray!55,
  fonttitle=\bfseries\small, title={#1}, boxrule=0.7pt, arc=2pt,
  left=6pt,right=6pt,top=4pt,bottom=4pt}

\newtcolorbox{answerbox}[1][Expected answer]{
  breakable, enhanced, colback=green!4, colframe=green!35!black,
  fonttitle=\bfseries\small, title={#1}, boxrule=0.7pt, arc=2pt,
  left=6pt,right=6pt,top=4pt,bottom=4pt}


\begin{document}


\LAtitle{Lab Assignment 3: Time-Series Regression: Stock Price Prediction}{Experiment 4 (Unit V) \textbullet\ CO2, CO4}{announced in the lab}{LQ3, after submission}

\section*{Objective}

\begin{itemize}[leftmargin=*]

  \item Build supervised features from an ordered series without leaking the future.

  \item Split time-ordered data correctly and see why a random split is wrong.

  \item Compare a model against the naive baseline that actually matters.

\end{itemize}

\section*{Datasets used in this assignment}

\textbf{C1 (daily stock OHLCV).}\\ Full descriptions, sources and loading snippets for all anchor datasets are in \texttt{Anchor-Datasets.pdf}, alongside this brief.

\section*{Instructions}

\begin{itemize}[leftmargin=*]

  \item Use a train/test split (80/20 unless stated) and set \texttt{random\_state} for reproducibility.

  \item Fit every transformer inside a \texttt{Pipeline} so nothing is fitted on the test set.

  \item Report the metric named in the question. A number without units or context earns no marks.

  \item Every question asks for a short written observation. Code without commentary caps you at half marks.

  \item The lab quiz that follows will ask you to read, trace and modify \emph{your own} submitted code.

\end{itemize}

\textbf{Submit:}

\begin{itemize}[leftmargin=*]

  \item Notebook: \texttt{AML\_LA03\_<SAPID>.ipynb} with all cells executed and outputs visible.

  \item Report: 2--3 pages PDF, or a \texttt{README.md}, carrying your results table and observations.

  \item Do not upload datasets. Cite the source and include the loading code.

\end{itemize}

\begin{center}\begin{tabular}{@{}llll@{}}\toprule
\textbf{Section} & \textbf{Level} & \textbf{Choice rule} & \textbf{Marks}\\
\midrule
A & Easy & Attempt any 2 of 3 (2 marks each) & 4\\
B & Medium & Attempt any 1 of 2 (4 marks each) & 4\\
C & Hard & Compulsory & 2\\
\midrule
\multicolumn{3}{@{}l}{\textbf{Total}} & \textbf{10}\\
\bottomrule\end{tabular}\end{center}

\noindent\small Course outcomes assessed by this assignment: \textbf{CO2, CO4}. Attempting more than the choice rule allows is not penalised; the best attempts are marked.\normalsize

\section*{Section A (Easy) --- attempt any 2 of 3 \quad (2 $\times$ 2 = 4)}

\begin{enumerate}[label=\textbf{A\arabic*.}, leftmargin=*]

  \item \textbf{Load and describe the series} \QMarks{2}\\
    \emph{Dataset: C1, one ticker, $\geq 3$ years.}
    \begin{enumerate}[label=(\alph*), leftmargin=*]
      \item Load the daily history, parse the date column and sort ascending.
      \item Plot the closing price and state the date range and number of trading days.
      \item State in two lines why the row order matters here but did not in LA02.
    \end{enumerate}

  \item \textbf{Lag features} \QMarks{2}\\
    \emph{Dataset: C1.}
    \begin{enumerate}[label=(\alph*), leftmargin=*]
      \item Create features \texttt{Close\_lag1}, \texttt{Close\_lag2}, \texttt{Close\_lag3} and a 5-day rolling mean.
      \item Set the target to the next day's close. Show the first five rows of the resulting frame.
      \item Explain in two lines why the first few rows must be dropped.
    \end{enumerate}

  \item \textbf{The naive baseline} \QMarks{2}\\
    \emph{Dataset: C1.}
    \begin{enumerate}[label=(\alph*), leftmargin=*]
      \item Build the baseline that predicts tomorrow's close equals today's close.
      \item Report its MAE and RMSE on your test period.
      \item State in two lines why any model must beat this number to be interesting.
    \end{enumerate}

\end{enumerate}

\section*{Section B (Medium) --- attempt any 1 of 2 \quad (1 $\times$ 4 = 4)}

\begin{enumerate}[label=\textbf{B\arabic*.}, leftmargin=*]

  \item \textbf{Random split versus chronological split} \QMarks{4}\\
    \emph{Dataset: C1.}
    \begin{enumerate}[label=(\alph*), leftmargin=*]
      \item Train a \texttt{RandomForestRegressor} twice: once with \texttt{train\_test\_split(shuffle=True)}, once with the last 20\% of dates held out.
      \item Report test MAE for both and put them side by side.
      \item Explain in 6--8 lines why the shuffled number is meaningless, naming the mechanism.
    \end{enumerate}

  \item \textbf{Horizon and difficulty} \QMarks{4}\\
    \emph{Dataset: C1, chronological split throughout.}
    \begin{enumerate}[label=(\alph*), leftmargin=*]
      \item Retrain your model to predict the close 1, 5 and 10 trading days ahead.
      \item Report MAE for each horizon alongside the persistence baseline at the same horizon.
      \item Plot MAE against horizon and explain the trend in 4--6 lines.
    \end{enumerate}

\end{enumerate}

\section*{Section C (Hard) --- compulsory \quad (2)}

\begin{enumerate}[label=\textbf{C\arabic*.}, leftmargin=*]

  \item \textbf{Would you trade on it?} \QMarks{2}\\
    \emph{Dataset: C1.}
    \begin{enumerate}[label=(\alph*), leftmargin=*]
      \item Convert your best model's predictions into an up/down direction call and report directional accuracy.
      \item In 4--6 lines, state whether a low MAE with near-50\% directional accuracy is useful, and why.
    \end{enumerate}

\end{enumerate}

\vfill\noindent\textit{End of Lab Assignment 3}

\notesources{\AMLLabSources}

\end{document}